\documentclass[11pt]{article}

% ACL style (review submission)
\usepackage[review]{acl}

% Standard package includes
\usepackage{times}
\usepackage{latexsym}

% For proper rendering and hyphenation of words containing Latin characters (including in bib files)
\usepackage[T1]{fontenc}
% For Vietnamese characters
% \usepackage[T5]{fontenc}
% See https://www.latex-project.org/help/documentation/encguide.pdf for other character sets

% This assumes your files are encoded as UTF8
\usepackage[utf8]{inputenc}

% This is not strictly necessary, and may be commented out,
% but it will improve the layout of the manuscript,
% and will typically save some space.
\usepackage{microtype}

% This is also not strictly necessary, and may be commented out.
% However, it will improve the aesthetics of text in
% the typewriter font.
\usepackage{inconsolata}

%Including images in your LaTeX document requires adding
%additional package(s)
\usepackage{graphicx}

% Additional packages for this paper
\usepackage{hyperref}
\usepackage{url}
\usepackage{booktabs}
\usepackage{amsfonts}
\usepackage{amsmath}
\usepackage{amssymb}
\usepackage{nicefrac}
\usepackage{xcolor}
\usepackage{multirow}
\usepackage{subcaption}
\usepackage{algorithm}
\usepackage{algorithmic}
\usepackage{tablefootnote}
\usepackage{pifont}
\newcommand{\cmark}{\ding{51}}
\newcommand{\xmark}{\ding{55}}

% If the title and author information does not fit in the area allocated, uncomment the following
%
%\setlength\titlebox{<dim>}
%
% and set <dim> to something 5cm or larger.
\setlength\titlebox{6cm}

\title{Structured Social State as an Inspectable Bottleneck for NPC Dialogue}

\author{Anonymous \\
  Affiliation / Address line 1 \\
  Affiliation / Address line 2 \\
  Affiliation / Address line 3 \\
  \texttt{email@domain} \\And
  Anonymous \\
  Affiliation / Address line 1 \\
  Affiliation / Address line 2 \\
  Affiliation / Address line 3 \\
  \texttt{email@domain} \\}

\begin{document}

\maketitle

\begin{abstract}
What information about an NPC's social state can be recovered from dialogue alone, and does that state support routing and constrained response generation?
We decompose NPC interaction into a 29-dimensional structured latent state $Z_t$ covering conversational intent, affect, relational stance, normative pressure, and decision policy.
A Qwen3-4B multi-head predictor shows uneven recoverability: arousal ($\kappa{=}0.692$), valence ($\kappa{=}0.688$), and response\_policy ($\kappa{=}0.632$) show the highest recoverability, while secrecy\_pressure ($\kappa{=}0.176$), familiarity\_delta ($\kappa{=}0.296$), and dominance\_delta ($\kappa{=}0.368$) are weak.
Other deltas retain measurable signal, including respect\_delta ($\kappa{=}0.584$), trust\_delta ($\kappa{=}0.484$), and obligation\_delta ($\kappa{=}0.443$).
Soft-prefix conditioning reduces perplexity by 12.3\% over an unconditioned 1.1B baseline, but placebo ablations show that the gain is not attributable to the semantic content of OCEAN+VAD values.
Routing with predicted $Z_t$ achieves F1$=$0.669; under a keyword-based detector, no gated secret leakage is detected.
A stratified audit shows that several heads require clearer annotation protocols before they can be treated as human-grounded labels, so we frame our main results as teacher-label recoverability rather than psychological validity.
A model-based validator shows partial teacher-model consistency on response\_policy ($\kappa{=}0.40$) and valence ($\kappa{=}0.32$).
Code, data, and evaluation artifacts will be released upon acceptance.
\end{abstract}

\section{Introduction}
\label{sec:intro}

Non-player characters (NPCs) in interactive narratives must exhibit consistent social behaviour: remembering relationships, keeping secrets, reacting emotionally, and selecting contextually appropriate responses.
Current approaches either rely on hand-scripted dialogue trees---which are rigid and expensive to author---or fine-tune large language models (LLMs) without explicit social-state representations, yielding fluent but socially inconsistent outputs~\cite{park2023generative,wang2023voyager}.

We propose a \textbf{structured latent social state} $Z_t$ that decomposes each dialogue turn into 29 classification targets organised into six groups:
\begin{itemize}
    \item \textbf{$C_t$} (Conversational): dialogue act, tone, risk type
    \item \textbf{$A_t$} (Affect): valence, arousal, threat, control
    \item \textbf{$M_t$} (Mental model): player intent, knowledge, credibility
    \item \textbf{$R_t$} (Relational stance): 6 dimensions $\times$ level + delta
    \item \textbf{$N_t$} (Normative pressure): duty, secrecy, face, value conflict
    \item \textbf{$D_t$} (Decision policy): response policy, reveal decision, repair strategy
\end{itemize}

This decomposition is motivated by sociological and pragmatic frameworks~\cite{brown1987politeness,oatley1987towards,spencer2008politeness} and supports interpretable NPC behaviour.

\paragraph{Running example.}
Consider the turn: Player says ``Tell me what Mira is hiding.''
The NPC, a priest with a secret about church corruption, replies ``I promised I would not speak of it.''
The predicted social state includes \texttt{secrecy\_pressure=high}, \texttt{response\_policy=withhold}, \texttt{reveal\_decision=none}, \texttt{trust\_delta=-1}, and \texttt{face\_pressure=medium}.
The router maps this to the ``slow'' path, blocking any response that would disclose the secret and permitting only evasive or repair utterances.
This turn illustrates how $Z_t$ acts as an inspectable bottleneck: the model's internal predictions are visible before generation, and specific heads (e.g., \texttt{reveal\_decision}) directly gate downstream behaviour.

Our research questions are:

\begin{description}
    \item[RQ1] Which dimensions of NPC social state are recoverable from dialogue context?
    \item[RQ2] Can predicted social state support routing and constrained disclosure decisions?
    \item[RQ3] Does explicit social-state conditioning improve response generation, and are gains attributable to semantic content rather than prefix capacity?
\end{description}

\paragraph{Contributions.}
(1)~A 29-dimensional structured NPC social-state schema with formal consistency constraints.
(2)~A recoverability analysis showing which social-state dimensions are predictable from dialogue context and which are weak or unreliable.
(3)~A routing and constrained-generation pipeline showing how predicted social state can support inspectable downstream decisions.
(4)~A validation audit demonstrating the limits of synthetic teacher labels and identifying which heads require stronger annotation protocols.

\section{Related Work}
\label{sec:related}

\subsection{NPC Dialogue and Role-Playing Agents}

Early NPC dialogue systems relied on finite-state machines and hand-authored dialogue trees~\cite{walker2006individual}.
Neural dialogue models replaced hand-written rules with sequence-to-sequence generation~\cite{bordes2017learning}, persona-conditioned dialogue~\cite{zhang2018personalizing}, and character-and-setting-conditioned generation in LIGHT~\cite{urbanek2019learning}.
More recent LLM-based agents extend this line to open-world and role-playing settings, including Generative Agents~\cite{park2023generative}, Voyager~\cite{wang2023voyager}, Character-LLM~\cite{shao2023character}, RoleLLM~\cite{wang2024rolellm}, and ChatHaruhi~\cite{li2023chatharuhi}.
These systems usually represent characters through free-text persona descriptions, prompts, or retrieved memories.
They do not expose a compact per-turn state that can be predicted, inspected, and used for routing.

Recent work has begun to study LLM-driven NPC dialogue directly~\cite{rimer2025talking,song2025llmnpc,buakhaw2025deflanderization,ploug2025open}.
These studies show the practical relevance of LLM-based NPCs, but most still rely on prompt-level control rather than a structured intermediate representation.
Our work instead evaluates a 29-head NPC-oriented state schema and measures which heads can be recovered from dialogue context.

\subsection{Structured Social State and Dialogue Tracking}

The proposed schema draws on prior work in affect, personality, politeness, and social reasoning, including the circumplex model of affect~\cite{russell1980circumplex}, cognitive theories of emotion~\cite{oatley1987towards}, the Big Five personality model~\cite{digman1990personality}, and politeness theory~\cite{brown1987politeness}.
Dialogue systems have used these signals for emotion-conditioned response generation~\cite{zhou2018emotional}, persona-based conversation~\cite{li2016persona}, and social simulation~\cite{gunkel2011computational}.
Most prior systems model one or two social dimensions rather than a multi-head social-state schema with consistency constraints.

Our formulation is also related to dialogue state tracking~\cite{mrksic2017neural,feng2023dialogue,jacqmin2024dialoguestate}, where a structured belief state is updated across turns.
Recent social-agent work has introduced explicit social representations, including SOTOPIA~\cite{zhou2024sotopia}, Social World Models~\cite{zhou2025social}, and future-aware state-action chains in SAGE~\cite{sage2025steering}.
Our contribution is narrower: we study an NPC-specific schema and evaluate recoverability, conditioning, and routing under a single synthetic dialogue pipeline.

\subsection{Conditioning and Controllable Generation}

Controllable generation methods include discrete control codes~\cite{keskar2019ctrl}, prefix tuning~\cite{li2021prefix}, soft prompting~\cite{lester2021power}, PPLM~\cite{dathathri2020plug}, and parameter-efficient tuning methods such as LoRA~\cite{hu2021lora} and QLoRA~\cite{dettmers2023qlora}.
These methods have been applied to sentiment, topic, and style control, but less often to multi-dimensional social state.

We use soft-prefix conditioning as a baseline and include placebo controls to test whether OCEAN+VAD values provide semantic control.
The placebo result motivates the structured $Z_t$ predictor: the prefix mechanism improves perplexity, but OCEAN+VAD values do not explain the gain.

\section{Method}
\label{sec:method}

\subsection{Social State Schema ($Z_t$)}
\label{sec:schema}

We define $Z_t$ as an operational control and inspection schema for NPC dialogue, not as a claim about human psychology. It is a teacher-labelled annotation schema over dialogue context and response policy; we do not treat it as psychological ground truth. Each dialogue turn $t$ is annotated with a structured social state $Z_t = (C_t, A_t, M_t, R_t, N_t, D_t)$, comprising 29 classification targets.
Table~\ref{tab:schema} summarises the schema.

\begin{table}[t]
\centering
\caption{Structured social state schema $Z_t$. Each field is a classification target predicted by the 29-head model. Stance dimensions ($R_t$) each have a level (5-class ordinal) and a per-turn delta (5-class ordinal).}
\label{tab:schema}
\small
\resizebox{\linewidth}{!}{%
\begin{tabular}{@{}llrl@{}}
\toprule
\textbf{Group} & \textbf{Field} & \textbf{\#Classes} & \textbf{Type} \\
\midrule
$C_t$ & dialogue\_act & 10 & multi-label \\
      & tone & 6 & single-label \\
      & risk\_type & 5 & single-label \\
\midrule
$A_t$ & valence, arousal, threat, control & 3 each & single-label \\
\midrule
$M_t$ & player\_intent & 9 & single-label \\
      & player\_knowledge & 4 & single-label \\
      & player\_credibility & 3 & single-label \\
\midrule
$R_t$ & \{affection, respect, dominance, & 5 each & ordinal \\
      & \ familiarity, trust, obligation\}$_\text{level}$ & & \\
      & \{...\}$_\text{delta}$ & 5 each & ordinal \\
\midrule
$N_t$ & duty, secrecy, face pressure & 3 each & single-label \\
      & value\_conflict & 3 & single-label \\
\midrule
$D_t$ & response\_policy & 10 & single-label \\
      & reveal\_decision & 4 & single-label \\
      & repair\_strategy & 5 & single-label \\
\bottomrule
\end{tabular}}
\end{table}

\paragraph{Consistency rules.}
We enforce domain constraints: e.g., $\text{reveal\_decision} = \textit{full}$ requires $\text{secrecy\_pressure} \neq \textit{high}$; $\text{response\_policy} = \textit{answer}$ implies $\text{reveal\_decision} \in \{\textit{hint}, \textit{partial}, \textit{full}\}$.
A consistency loss (Section~\ref{sec:joint}) penalises violations during joint training.

\subsection{Architecture Stack}
\label{sec:architectures}

We organise our benchmark into four tracks (Appendix~Figure~\ref{fig:architecture_stack}):

\paragraph{Track A: From-scratch SLMs (15--22M).}
Four architectures share the same vocabulary (GPT-2 tokeniser, 50,257 tokens), embedding dimension ($d=512$), and training data (16,905 dialogue lines):

\begin{enumerate}
    \item \textbf{GPT}: 6-layer transformer decoder, 8 attention heads. Parameters: 16.1M.
    \item \textbf{PrefixGPT}: GPT + learned soft-prefix from an 8D conditioning vector (5 OCEAN + 3 VAD), projected to 4 prefix tokens via a 2-layer MLP. Parameters: 16.6M.
    \item \textbf{MoE}: GPT with Mixture-of-Experts FFN layers (4 experts, top-2 routing with load-balancing loss). Each expert adds separate FFN parameters, yielding 22.4M total.
    \item \textbf{Mamba-like}: Simplified state-space model (SSM) with selective scan, implemented in pure PyTorch. Parameters: 15.4M.
\end{enumerate}

\paragraph{Track B: Conditioning encoders.}
DistilBERT-based encoders predict personality (OCEAN, 5D regression) and affect (VAD, 3D regression) from dialogue context.
These produce the conditioning vector $c_t \in \mathbb{R}^8$ used by PrefixGPT and ConditionalDialogue.

\paragraph{Track C: Response generation.}
Three models generate NPC responses:
\begin{enumerate}
    \item \textbf{ConditionalDialogue}: From-scratch GPT + personality/affect soft-prefix embeddings ($c_t \to$ prefix via learned projection).
    \item \textbf{TinyLlama + LoRA}: 1.1B pretrained chat model, fine-tuned with LoRA ($r=16$, $\alpha=32$) on SFT data without social-state conditioning.

    \item \textbf{Gemma + QLoRA}: exploratory larger pretrained baselines, including a corrected Gemma-2-2B-it run and a Gemma-4-E2B engineering run.
\end{enumerate}

\paragraph{Track D: Structured LLM.}
A three-stage pipeline on Qwen3:
\begin{enumerate}
    \item \textbf{Latent predictor}: Qwen3-4B backbone + LoRA + 29 classification heads. Each head is a 2-layer MLP ($\text{hidden} \to 256 \to n_\text{classes}$) applied to the last-token hidden state.
    \item \textbf{Response generator}: Qwen3-4B + LoRA ($r=32$), supervised on input $\to$ response with gold $Z_t$ in the prompt.
    \item \textbf{Joint model}: Shared backbone with both classification heads and LM loss, plus a consistency loss $\mathcal{L}_\text{consist}$ penalising $P(\text{full\_reveal}) \cdot P(\text{high\_secrecy})$.
\end{enumerate}

\subsection{Training}
\label{sec:training}

\paragraph{Loss functions.}
Track A uses standard cross-entropy LM loss.
Track B uses MSE for regression targets.
Track D Stage~1 uses a weighted multi-head loss:
\begin{equation}
    \mathcal{L}_\text{heads} = \sum_{g \in \{C,A,M,R,N,D\}} \lambda_g \cdot \frac{1}{|g|} \sum_{f \in g} \mathcal{L}_\text{CE}(f)
\end{equation}
with $\lambda_R = \lambda_D = 2.0$, $\lambda_M = 1.5$, others $= 1.0$, reflecting the importance of relational and decision heads.

\paragraph{Joint loss (Stage 3).}
\label{sec:joint}
\begin{equation}
    \begin{aligned}
    \mathcal{L}_\text{joint}
    &= \mathcal{L}_\text{heads} + \lambda_Y \mathcal{L}_\text{LM} \\
    &\quad + \lambda_\text{consist}\,\mathbb{E}[p_\text{full}p_\text{secret}].
    \end{aligned}
\end{equation}
Here $p_\text{full}=P(\text{full\_reveal})$ and $p_\text{secret}=P(\text{high\_secrecy})$.

\paragraph{Optimisation.}
All models use AdamW.
Track A: $\text{lr}=3 \times 10^{-4}$, batch 32, 20 epochs.
Track D: $\text{lr}=2 \times 10^{-4}$ (backbone), $4 \times 10^{-4}$ (heads), cosine schedule, gradient accumulation $=32$, BFloat16.

The full data flow from scenario generation through validation and training is shown in Appendix~Figure~\ref{fig:data_flow}.

\subsection{Fast/Slow Routing}
\label{sec:routing}

A rule-based router classifies turns into \textit{fast path} (direct generation) or \textit{slow path} (full $Z_t$ inference) based on predicted $D_t$ and $N_t$.
The slow path is used when value conflict is strong, secrecy is high with a non-\texttt{none} reveal decision, or the response policy is \texttt{threaten} or \texttt{negotiate}.

\section{Experimental Setup}
\label{sec:experiments}

\paragraph{Data.}
We use 7,742 dialogue turns across 736 episodes spanning 7 scenario types (secret extraction, trust building, alliance negotiation, threat escalation, deception detection, rumor confrontation, apology repair).
Each turn is annotated with the full 29-field $Z_t$ schema (\texttt{train\_heads.jsonl}, \texttt{val\_heads.jsonl}, \texttt{test\_heads.jsonl}).
Data is split 80/10/10 into train (6,175 turns, 587 episodes), validation (683 turns, 69 episodes), and test (884 turns, 80 episodes).
Appendix~Table~\ref{tab:scenario_dist} shows the scenario-type distribution across splits; all seven scenario types are represented in each split, with secret extraction and rumor confrontation being the most frequent.
The mean turns per episode is 10.5 (train), 9.9 (val), and 11.1 (test); mean context length is 200, 190, and 219 words respectively.
SLM training (Track~A) uses a separate corpus of 16,905 plain-text dialogue lines (2,183 structured JSON records after blank-line segmentation) spanning 396 unique NPCs (\texttt{train.txt} / \texttt{val.txt} split at 95/5).
The SLM corpus covers all 7 scenario types, with secret extraction (20.1\%), rumor confrontation (18.4\%), and apology repair (17.2\%) being the most frequent; response lengths average 29.8 words and context lengths average 141 words (Table~\ref{tab:slm_dist} in Appendix~\ref{app:datagen}).
Encoder training (Track~B) uses 2,221 personality-labelled samples from essays-big5\footnote{\url{https://huggingface.co/datasets/jingjietan/essays-big5}} and 9,056 affect-labelled samples drawn from EmoBank~\cite{buechel2017emobank} and EmpatheticDialogues~\cite{rashkin2019empathetic}.
\textbf{Data provenance:} Source documentation, including which datasets are real (published corpora, human-annotated) versus synthetic (teacher LLM-generated), with sizes and merge details, is provided in Appendix~\ref{app:datagen}.
Trained model checkpoints, generated data splits, scenario definitions, evaluation code, and compact evaluation artifacts will be released upon acceptance.

\paragraph{Hardware.}
Appendix~Table~\ref{tab:hardware} summarises the memory requirements for training and inference across all models.
Reported values include model weights, KV-cache for inference (sequence length 256 for Tracks~A/B, 2048 for Tracks~C/D), and optimizer states for training (AdamW: 2$\times$ fp32 copies of parameters).

All training used NVIDIA A30 GPUs (24\,GB), CUDA 12.4, PyTorch 2.6.
The Gemma\,4\,E2B model required the PyTorch \texttt{expandable\_segments} CUDA allocator during evaluation only (activation memory spike); training fits without memory flags.
Tracks\,A and~B train on any GPU with 4\,GB or more VRAM, or on CPU alone with longer training times.

\paragraph{Metrics.}
\textit{Track A}: validation perplexity (PPL).
\textit{Track B}: F1 (personality), CCC (affect), MSE.
\textit{Track C}: PPL, ROUGE-L, BLEU-1/2/4, Distinct-1/2/3, secret leakage rate, contradiction rate.
\textit{Track D}: per-head accuracy, macro-F1, true Cohen's $\kappa$ computed from confusion matrices, per-group means, response policy F1.
\textit{Routing}: precision, recall, F1, false positive rate.

\paragraph{Statistical reporting.} Unless otherwise noted, all mean perplexity values are averaged over independent random seeds with standard deviation reported.
Bootstrap 95\% confidence intervals ($n{=}1{,}000$ resamples) are reported for ROUGE-L and paired comparisons where indicated.
Single-seed results are explicitly labelled as such.

\section{Results}
\label{sec:results}

\subsection{From-Scratch SLM Baseline}

\begin{table}[t]
\centering
\caption{From-scratch SLM validation perplexity (20 epochs, 3 seeds). \textbf{Bold}: best mean. GPT-22M is param-matched to MoE at $\sim$22M.}
\label{tab:slm}
\small
\resizebox{\linewidth}{!}{%
\begin{tabular}{@{}lrrrrr@{}}
\toprule
\textbf{Architecture} & \textbf{Params (M)} & \textbf{Mean val\_ppl $\downarrow$} & \textbf{Std} & \textbf{$\Delta$ vs GPT-16M} & \textbf{Conditioning} \\
\midrule
\textbf{GPT-22M} & 22.3 & \textbf{41.65} & 0.88 & $-6.1\%$ & none \\
MoE       & 22.4 & 43.38 & 0.91 & $-2.2\%$ & none \\
GPT-45M (seed 42)   & 44.8 & 41.61 & --- & $-6.2\%$ & none \\
PrefixGPT          & 16.6 & 44.11          & 0.63 & $-0.6\%$ & OCEAN+VAD \\
GPT-16M                & 16.1 & 44.37          & 0.58 & ---      & none \\
Mamba-like         & 15.4 & 53.45          & 0.35 & $+20.5\%$& none \\
\bottomrule
\end{tabular}}
\end{table}

When controlling for parameter budget, a dense GPT at 22.3M parameters achieves the best mean perplexity (41.65 across 3 seeds), outperforming MoE at 22.4M (43.38) by 4.0\% and the smaller GPT-16M (44.37) by 6.1\%.
This reverses the naive comparison that suggested MoE was architecturally superior: the apparent MoE advantage over GPT-16M was a parameter-count artifact, not an expert-routing benefit.
Scaling a standard dense GPT from 16.1M to 22.3M parameters yields a larger perplexity reduction than adding expert routing and load-balancing at the same total parameter count.
Further scaling to 44.8M (GPT-45M, single seed) produces only a marginal additional gain (41.61 vs. 41.65), suggesting that the $\sim$22M range is near the point of diminishing returns for this dataset size (16,905 dialogue lines).
PrefixGPT's modest gain ($-0.6\%$) over GPT-16M with only a 0.5M parameter increase suggests the 8D conditioning signal provides marginal but consistent benefit at this scale.
The Mamba-like SSM converges fastest (best loss at epoch 10) but plateaus at a higher perplexity, consistent with prior observations that SSMs can underperform transformers on tasks requiring long-range attention over structured input~\cite{gu2023mamba}.

\textbf{Metric caveat.} Cross-architecture perplexity is reported descriptively; because tokenization differs across model families, like-for-like comparisons rely on matched settings or normalized alternatives such as bits-per-byte.
Within the GPT family (GPT-16M, GPT-22M, GPT-45M), tokenizer identity ensures direct comparability.

\subsection{Track B: Conditioning Encoders}

\begin{table}[t]
\centering
\caption{Conditioning encoder evaluation.}
\label{tab:encoders}
\small
\resizebox{\linewidth}{!}{%
\begin{tabular}{@{}lrrrrr@{}}
\toprule
\textbf{Encoder} & \textbf{Base} & \textbf{val\_F1 $\uparrow$} & \textbf{val\_CCC $\uparrow$} & \textbf{val\_MSE $\downarrow$} & \textbf{Best Epoch} \\
\midrule
Personality (OCEAN) & DistilBERT & 0.678 & --- & 0.248 & 4/15 \\
Affect (VAD)        & DistilBERT & ---   & 0.559 & 0.005 & 13/15 \\
\bottomrule
\end{tabular}}
\end{table}

Personality prediction achieves F1~$=0.678$, consistent with known difficulty of text$\to$OCEAN regression~\cite{mairesse2007using}.
The affect encoder achieves CCC~$=0.559$, supporting valence-arousal-dominance tracking.
A cache of 414 NPC personality embeddings was built from training data for efficient inference.

\subsection{Response Conditioning (RQ3)}

\begin{table}[t]
\centering
\caption{Response generation comparison. \textbf{Bold}: best PPL. The soft-prefix gain is $12.3\%$, but placebo ablations show this gain is not attributable to the semantic content of OCEAN+VAD values. ConditionalDialogue reports mean across 3 seeds (42--44).}
\label{tab:response_ppl}
\small
\resizebox{\linewidth}{!}{%
\begin{tabular}{@{}llrrr@{}}
\toprule
\textbf{Model} & \textbf{Conditioning} & \textbf{val\_ppl $\downarrow$} & \textbf{Epochs} & \textbf{Date} \\
\midrule
\textbf{ConditionalDialogue} (3 seeds) & OCEAN+VAD soft-prefix \hfill & \textbf{2.88} (mean) & 5 & May 4--5 \\
TinyLlama 1.1B + LoRA       & none (SFT)            & 3.30          & 3 & May 2 \\
Gemma-2-2B-it + QLoRA      & NPC profile (SFT)     & 6.38          & 2 & May 3 \\
Gemma-4-E2B + QLoRA          & none (SFT)            & 13.48         & 3 & May 20 \\
Gemma-4-E2B + social-state   & gold $Z_t$ XML prefix  & 13.61         & 3 & May 20 \\
\bottomrule
\end{tabular}}
\end{table}

These PPL values are used as within-experiment diagnostics and are not tokenizer-normalized across model families.

For RQ3, soft-prefix conditioning reduces response perplexity by \textbf{12.3\%} (3.30$\to$2.88 mean across three seeds; Table~\ref{tab:response_ppl}).
However, placebo controls with shuffled OCEAN values, random VAD values, and shuffled+random conditioning all converge to the same PPL range (2.88--2.91), indicating that the gain comes from the presence of prefix-conditioning capacity rather than from the semantic content of OCEAN/VAD vectors.
A multi-seed placebo experiment (3 seeds $\times$ real vs. random VAD; Table~\ref{tab:placebo}) is consistent with this null result: all six runs converge to val\_ppl $\in [2.859, 2.900]$, a spread of only $1.4\%$ with no systematic difference between real and randomised VAD conditioning.
This is still a useful architectural result, but it rules out the stronger claim that the learned OCEAN/VAD values causally improve generation.
The corrected Gemma-2-2B-it run (PPL 6.38) provides the primary pretrained baseline, while the Gemma-4-E2B reruns at 3 epochs (baseline PPL 13.48) show that gold $Z_t$ XML prefix conditioning provides zero improvement (13.48$\to$13.61), indicating that this prompt format does not improve over simple SFT at this model scale.

The placebo table and response-model chart are provided in Appendix~Table~\ref{tab:placebo} and Appendix~Figure~\ref{fig:response_ppl}.

\begin{table}[t]
\centering
\caption{Response generation quality metrics (Qwen3-4B response model, $n=683$ validation turns).}
\label{tab:response_quality}
\small
\begin{tabular}{@{}lr@{}}
\toprule
\textbf{Metric} & \textbf{Value} \\
\midrule
ROUGE-L (bootstrap 95\% CI) & 0.145 [0.140, 0.151] \\
BLEU-1 / BLEU-4 & 0.273 / 0.049 \\
Distinct-2 & 0.638 \\
Detected gated leakage (reveal=none) & 0.000 \\
Ungated secret leakage rate & 0.076 \\
Contradiction rate & 0.000 \\
Prompt artifact rate & 0.000 \\
Avg generation length / ref length & 29.3 / 32.2 \\
Length ratio (gen/ref) & 0.909 \\
Repeated 3-gram rate & 0.000 \\
\bottomrule
\end{tabular}
\end{table}

Table~\ref{tab:response_quality} reports the final automatic response evaluation on the Qwen3-4B response generator.
ROUGE-L of 0.145 exceeds the unconditioned baseline (ROUGE-L$=$0.115), with properly controlled generation length (ratio$=$0.91) and high lexical diversity (Distinct-2$=$0.638).
Detected gated secret leakage---the safety-relevant metric, computed only over turns where \texttt{reveal\_decision=none}---is zero under our keyword-based heuristic detector.
Ungated leakage across all 683 turns is 7.6\%, reflecting that the keyword-based detector flags some non-secret patterns; production reliability would require manual verification or a trained classifier.
Contradictions and prompt artifacts are both zero, consistent with the interpretation that the length-aware decoding configuration resolves the earlier verbosity and formatting issues.

Test-set evaluation on 884 held-out turns (Table~\ref{tab:test_set}) shows similar test behaviour with modest degradation: ROUGE-L drops from 0.145 to 0.139 ($-4.3\%$), Distinct-2 from 0.638 to 0.603, and length ratio from 0.91 to 0.86. Detected gated leakage remains zero on test under the same heuristic detector.
These automatic metrics are used as development diagnostics and do not measure player-perceived dialogue quality.

\begin{table}[t]
\centering
\caption{Validation vs. test-set comparison (Qwen3-4B pipeline).}
\label{tab:test_set}
\small
\resizebox{\linewidth}{!}{%
\begin{tabular}{@{}lrrr@{}}
\toprule
\textbf{Metric} & \textbf{Val (683)} & \textbf{Test (884)} & \textbf{$\Delta$} \\
\midrule
\textbf{Response Generation} & & & \\
ROUGE-L & 0.145 & 0.139 & $-4.3\%$ \\
BLEU-4 & 0.049 & 0.039 & $-20\%$ \\
Distinct-2 & 0.638 & 0.603 & $-5.5\%$ \\
Length ratio & 0.909 & 0.864 & $-5.0\%$ \\
Detected gated leakage & 0.000 & 0.000 & = \\
Ungated leakage & 0.076 & 0.067 & $-0.9$pp \\
\midrule
\textbf{Latent Prediction} & & & \\
Mean accuracy & 0.688 & 0.673 & $-2.1\%$ \\
Mean Cohen's $\kappa$ & 0.484 & 0.461 & $-4.8\%$ \\
Response-policy F1 & 0.622 & 0.427 & $-31\%$ \\
Reveal-decision F1 & 0.656 & 0.595 & $-9.3\%$ \\
\midrule
\textbf{Routing (predicted $Z_t$)} & & & \\
Routing F1 & 0.669 & 0.686 & +2.5\% \\
False positive rate & 0.399 & 0.318 & $-20\%$ \\
Slow-path rate & 0.548 & 0.501 & $-8.6\%$ \\
\bottomrule
\end{tabular}}
\end{table}

\subsection{Latent State Prediction (RQ1)}

\begin{table}[t]
\centering
\caption{Per-group latent state prediction metrics (Qwen3-4B, 28 single-label heads evaluated on validation set; dialogue act is multi-label, excluded here). All $\kappa$ values are true Cohen's $\kappa$ computed from confusion matrices.}
\label{tab:latent_groups}
\small
\resizebox{\linewidth}{!}{%
\begin{tabular}{@{}llrrrr@{}}
\toprule
\textbf{Group} & \textbf{Description} & \textbf{\#Heads} & \textbf{Mean Acc $\uparrow$} & \textbf{Mean Macro-F1} & \textbf{Mean $\kappa$ $\uparrow$} \\
\midrule
$A_t$ & Affect           & 4  & 0.794 & 0.714 & 0.645 \\
$N_t$ & Normative        & 4  & 0.727 & 0.597 & 0.364 \\
$M_t$ & Mental model     & 3  & 0.726 & 0.521 & 0.487 \\
$D_t$ & Decision policy  & 3  & 0.676 & 0.586 & 0.508 \\
$C_t$ & Conversational   & 2  & 0.679 & 0.441 & 0.469 \\
$R_t$ & Relational stance& 12 & 0.636 & 0.483 & 0.467 \\
\midrule
\multicolumn{2}{l}{\textbf{Overall (28 heads)}} & 28 & \textbf{0.688} & \textbf{0.541} & \textbf{0.484} \\
\bottomrule
\end{tabular}}
\end{table}

For RQ1, the 29-head predictor on Qwen3-4B reveals a systematic pattern in what is---and is not---learnable from dialogue context (Table~\ref{tab:latent_groups}).
Static attributes achieve higher $\kappa$ than most relational-delta heads: arousal (true $\kappa{=}0.692$), valence ($\kappa{=}0.688$), and duty\_pressure ($\kappa{=}0.629$) all show strong signal.
Per-turn \textit{deltas} show a mixed pattern: familiarity\_delta ($\kappa{=}0.296$) and dominance\_delta ($\kappa{=}0.368$) are weakly recoverable, while respect\_delta ($\kappa{=}0.584$), trust\_delta ($\kappa{=}0.484$), and obligation\_delta ($\kappa{=}0.443$) retain measurable signal.
This variation suggests that inferring fine-grained social change from a single turn of text is not uniformly difficult; some relational dimensions are recoverable while others are weak.
The actionable response\_policy head (10 classes) achieves macro-F1$=$0.622 (true $\kappa{=}0.632$), well above the 0.10 chance baseline.
The reveal\_decision head achieves macro-F1$=$0.656.
Overall mean accuracy is 0.688 with true Cohen's $\kappa{=}0.484$ across 28 evaluated heads.
This true $\kappa$ is a conservative metric; normalized above-chance estimates (which ignore class priors) would be $\sim$0.611.
Under the heuristic detector, secret leakage---defined as predicting a reveal decision other than \texttt{none} when secrecy\_pressure is \texttt{high}---is zero.
Appendix~\ref{app:prediction_examples} shows concrete prediction examples on held-out test turns.

\begin{figure}[t]
\centering
\includegraphics[width=0.86\linewidth]{figures/latent_group_scores.pdf}
\caption{Latent prediction by social-state group (Qwen3-4B, true Cohen's $\kappa$). Affect has the highest mean $\kappa$, while normative-pressure heads have the lowest mean $\kappa$ despite relatively high accuracy, indicating class-prior effects.}
\label{fig:latent_groups}
\end{figure}

Latent-predictor training dynamics are shown in Appendix~Figure~\ref{fig:qwen_training}.

\subsection{Routing and Constrained Generation (RQ2)}

The rule-based fast/slow router achieves F1$=$0.669 when evaluated with \textit{predicted} latent states from the Qwen3-4B predictor (Table~\ref{tab:routing}).
This is the routing-time metric: the router applies deterministic rules to the predicted $\hat{Z}_t$ rather than to gold labels.
Precision is 0.663, recall is 0.676, with a false positive rate of 0.399 and a slow-path rate of 54.8\% on the validation set.
Prediction coverage is 100\% (all 683 validation turns have a predicted $Z_t$).
On the test set (363 unique turns), routing achieves F1$=$0.686 (precision 0.648, recall 0.728, FPR 0.318), slightly better than validation, indicating that the predicted-state router does not degrade on this test subset.
Gold-state routing (a deterministic sanity check with oracle $Z_t$) achieves F1$=$1.0; the F1$\approx$0.67--0.69 values reported here are the real generalization test with predicted states.

\begin{table}[t]
\centering
\caption{Routing performance with predicted latent states (Qwen3-4B).}
\label{tab:routing}
\small
\begin{tabular}{@{}lr@{}}
\toprule
\textbf{Metric} & \textbf{Value} \\
\midrule
Routing F1 & 0.669 \\
Precision & 0.663 \\
Recall & 0.676 \\
False positive rate & 0.399 \\
Slow-path rate & 0.548 \\
Prediction coverage & 1.000 \\
\bottomrule
\end{tabular}
\end{table}

Joint training did not improve mean latent prediction over the separate-stage pipeline ($\Delta\kappa{=}-0.024$); full per-head results appear in Appendix~\ref{app:joint}.
Appendix~\ref{app:jepa} reports a mixed-to-null Social-State JEPA exploratory result.
Appendix~\ref{app:gold_zt_bridge} reports per-field gold-$Z_t$ response-difficulty analysis.

\section{Discussion}
\label{sec:discussion}

\paragraph{Audit as methodological contribution.}
The stratified audit does not invalidate the teacher-label recoverability experiment; rather, it indicates that the current annotation guidelines were insufficient for untrained annotators, not that the schema is unusable.
Stronger protocols---fewer heads per task, trained annotators, a calibration round, adjudication, and clearer definitions---would be needed before these heads can be treated as human-grounded labels.
Some low-agreement heads may simply need to be collapsed or redefined.
This turns the audit from a weakness into a serious methodological contribution: it identifies which dimensions of social state are annotation-sensitive and which are robust.

\paragraph{Architecture choice at small scale.}
At the 15--22M parameter scale, a dense GPT outperforms MoE when matched for parameter count (PPL 41.65 vs. 43.38), and pretraining dominates all architectural choices (best from-scratch PPL 41.65 vs. pretrained 2.88).

\paragraph{Conditioning effectiveness.}
The 12.3\% perplexity reduction from soft-prefix conditioning is notable because it requires no additional inference-time compute beyond the prefix projection.
However, placebo ablations show that shuffled OCEAN values and random VAD values achieve the same PPL range as real OCEAN+VAD conditioning.
Thus the improvement should be interpreted as a prefix-capacity or optimization effect, not as evidence that OCEAN/VAD values provide semantically useful control.
This negative result motivates structured $Z_t$ conditioning and Social-State JEPA as more targeted mechanisms for learning future-relevant social dynamics.

\paragraph{Latent state as an inspectable bottleneck.}
The 29-head predictor on Qwen3-4B indicates that parts of structured social state are learnable from dialogue context.
Mean true Cohen's $\kappa=0.484$ across 28 heads indicates substantial lift over chance, with the best groups (Affect $\kappa=0.645$, Decision $\kappa=0.508$) showing strong signal.
Individual heads such as \texttt{secrecy\_pressure=high} and \texttt{response\_policy=withhold} expose intermediate NPC state predictions before generation.
Response policy F1$=$0.622 and reveal decision F1$=$0.656 indicate that the most actionable downstream heads are among the best predicted.

\paragraph{Secret-keeping.}
Detected gated secret leakage---computed only over turns where \texttt{reveal\_decision=none}---is zero under our keyword-based heuristic detector across all 683 validation turns.
Ungated leakage across all turns is 7.6\%, reflecting keyword-based heuristic false positives that are outside the scope of this heuristic evaluation.
This zero-detected-leakage result supports the latent predictor's $D_t$ output as a promising mechanism for gating information disclosure, though production reliability requires a validated classifier or human audit.

\section{Conclusion}
\label{sec:conclusion}

We set out to determine whether structured social state can be recovered from NPC dialogue and whether it supports routing and constrained response generation.
Four findings emerge.

First, social state is partially recoverable. The recoverability pattern is systematic.
Static attributes---emotional arousal, valence, duty pressure---achieve higher $\kappa$ than most relational-delta heads (true $\kappa{>}0.62$).
Dynamic per-turn changes are not uniformly difficult: familiarity\_delta and dominance\_delta are weakly recoverable, while respect\_delta and trust\_delta retain measurable signal.
This variation suggests that single-turn social-state prediction has limits, but the difficulty differs across relational dimensions.

Second, soft-prefix conditioning reduces perplexity, but OCEAN/VAD values do not explain the gain.
The conditioned model reduces perplexity by 12.3\% over an unconditioned baseline trained on identical data, but shuffled and random conditioning controls match the same PPL range.
Gemma-4-E2B baselines with gold $Z_t$ XML prefix show zero improvement over unconditioned SFT (13.48 vs. 13.61), so this XML prefix format does not improve over unconditioned SFT at this model scale.

Third, the audit identifies which heads require stronger annotation protocols.
The audit flags that secrecy pressure, familiarity level, and some relational deltas need clearer guidelines or trained annotators before they can be treated as human-grounded labels.
This is a methodological contribution, not merely a limitation.

Fourth, the structured pipeline runs end-to-end with predicted states.
Routing with predicted $Z_t$ achieves F1$=$0.669; response generation with predicted states achieves ROUGE-L$=$0.145 with zero detected gated leakage under the keyword-based detector.
These numbers are modest but support the full predicted-state pipeline as functional, not merely a gold-label thought experiment.

Code, data, and evaluation artifacts will be released upon acceptance.
The structured latent state $Z_t$ provides an inspectable intermediate representation between perception and generation---one that is predictable from dialogue context, but the audit indicates that several heads require better annotation protocols before they can be treated as human-grounded labels.

\section*{Limitations}
\label{sec:limitations}

\paragraph{Synthetic labels and audit validity.}
All dialogue turns and social-state labels are generated by a teacher LLM pipeline.
The recoverability results therefore measure agreement with teacher-generated labels, not agreement with latent human social judgments.
High $\kappa$ indicates that the model can reproduce the teacher-labelled schema from dialogue context; it does not imply psychological validity.

We conducted a stratified audit of 150 natural, non-counterfactual test turns across eight heads: secrecy pressure, reveal decision, response policy, repair strategy, valence, arousal, trust level, and familiarity level.
Five uncalibrated human annotators and one model-based validator were used.
All five human annotators were flagged by post-hoc quality-control checks, and human--human agreement was low for most heads.
The model-based validator showed partial teacher-model consistency on \texttt{response\_policy} ($\kappa{=}0.40$) and \texttt{valence} ($\kappa{=}0.32$).
These results do not invalidate the teacher-label recoverability experiment, but they show that several heads require clearer annotation protocols or trained annotators before they can be treated as human-grounded labels.

\paragraph{Data scale and domain.}
The dataset contains 7,742 turns across 736 episodes in a single fantasy setting, ``Oakhaven Siege''.
Results may change with larger data, different genres, open-world dialogue, or non-English settings.
Rare decision classes are especially sensitive to data scale.

\paragraph{Automatic response metrics.}
BLEU, ROUGE-L, Distinct, and perplexity are weak proxies for dialogue quality, role consistency, and player-perceived usefulness.
We use them as development diagnostics.
Claims about NPC response quality require human preference evaluation or task-success measurement.

\paragraph{Heuristic leakage detection.}
The leakage detector is keyword-based and has not been validated for precision or recall.
``Zero detected gated leakage'' means zero detections under this heuristic, not zero leakage in principle.
A trained detector or manual verification is required before using the gating mechanism in production.

\paragraph{Model comparisons.}
Several comparisons are diagnostic rather than definitive.
Cross-family perplexity comparisons are affected by tokenizer and training differences.
The joint model shows a small mean degradation relative to the separate-stage pipeline ($\Delta\kappa{=}-0.024$), and per-head trade-offs require additional paired tests.

\section*{Ethical Considerations}

This work raises several ethical issues that deserve explicit discussion.

\textbf{Synthetic data provenance.} All dialogue and social-state labels were generated by a teacher LLM pipeline rather than collected from human players or annotators.
While this enables controlled experimentation at small scale, the labels reflect the teacher model's biases and may not generalise to real human social behaviour.
The generation prompts and consistency rules will be released upon acceptance so that others can audit the synthetic supervision chain.

\textbf{Stereotype and bias risk.} The 29-head schema includes personality (OCEAN), affect (VAD), and relational dimensions that could encode demographic stereotypes if mapped to NPC profiles without care.
Our synthetic profile generator samples from uniform distributions and does not condition on protected attributes, but downstream practitioners who reuse the schema for commercial games should validate that character traits do not correlate with harmful stereotypes.

\textbf{Deception and manipulation.} The secrecy-gating mechanism is designed to withhold information from a player character, which is a form of in-game deception.
While this is standard in narrative design, deploying such systems in open-world settings with real human interlocutors raises concerns about manipulative NPC behaviour.
We explicitly scope the intended use to fictional game environments with clear player consent.

\textbf{Privacy and consent.} The dataset contains no real human data; all turns are synthetic.
However, the scenario templates draw on common social situations (negotiation, confession, blackmail) that could be traumatic for some players.
Any deployment should include content warnings and opt-in mechanisms.

\textbf{Environmental cost.} Training the from-scratch SLMs required approximately 20 GPU-hours on an NVIDIA A30; the pretrained-model experiments required approximately 120 GPU-hours.
We report these budgets to support reproducibility and carbon-accounting efforts.

\bibliography{references}

\appendix

\section{Positioning and Pipeline Figures}
\label{app:positioning}

\begin{table}[tp]
\centering
\caption{How our structured social-state framework advances beyond prior work.}
\label{tab:novelty}
\footnotesize\setlength{\tabcolsep}{2pt}
\begin{tabular}{@{}p{0.22\linewidth}p{0.33\linewidth}p{0.35\linewidth}@{}}
\toprule
\textbf{Research Line} & \textbf{What Prior Work Does} & \textbf{What We Add} \\
\midrule
NPC dialogue~\cite{park2023generative,wang2023voyager,urbanek2019learning}
    & Generate dialogue from persona strings or without explicit state
    & 29-dim structured state $Z_t$ with formal consistency constraints; inspectable intermediate representation \\
\midrule
Social simulation~\cite{oatley1987towards,russell1980circumplex,digman1990personality}
    & Model one or two social dimensions (affect, personality) in isolation
    & Unify 6 social groups into a single predictable, verifiable, actionable state; cross-dimensional consistency rules \\
\midrule
Conditioned generation~\cite{li2021prefix,keskar2019ctrl,dathathri2020plug}
    & Control single attributes (sentiment, formality, topic)
    & Placebo-controlled multi-dimensional conditioning analysis; structured $Z_t$ state integrated into a 3-stage pipeline with joint training \\
\midrule
Multi-head classification~\cite{liu2019multi,li2021past,raheja2019dialogue}
    & Classify single-turn attributes (emotion, dialogue acts)
    & Multi-turn structured state prediction across 6 qualitatively distinct groups; feeds back into generation via joint training \\
\bottomrule
\end{tabular}
\end{table}

\begin{figure}[tp]
\centering
\includegraphics[width=0.92\linewidth]{figures/architecture_stack.pdf}
\caption{Four-track benchmark stack. The same scenario-generated data feeds architecture benchmarking, social-state encoders, response models, and the structured Qwen3 pipeline.}
\label{fig:architecture_stack}
\end{figure}

\begin{figure}[tp]
\centering
\includegraphics[width=\linewidth]{figures/data_flow.pdf}
\caption{Data flow from scenario specifications to validated splits and paper artifacts. Each generated turn is validated against the $Z_t$ schema before use in training or evaluation.}
\label{fig:data_flow}
\end{figure}

\begin{figure}[tp]
\centering
\includegraphics[width=0.96\linewidth]{figures/best_model_diagram.pdf}
\caption{ConditionalDialogue soft-prefix response model. The response model projects an 8D OCEAN+VAD vector into soft-prefix tokens before LoRA response generation; placebo ablations show the prefix mechanism helps, while the semantic content of OCEAN/VAD values does not.}
\label{fig:best_model}
\end{figure}

\begin{figure}[tp]
\centering
\includegraphics[width=0.82\linewidth]{figures/slm_ppl_comparison.pdf}
\caption{Track A validation perplexity (3 seeds, mean $\pm$ std). GPT-22M (22.3M) beats MoE (22.4M) by 4.0\% at equal parameter count, reversing the parameter-confounded comparison.}
\label{fig:slm_ppl}
\end{figure}

\begin{figure}[tp]
\centering
\includegraphics[width=0.88\linewidth]{figures/response_ppl_comparison.pdf}
\caption{Track C response-model validation perplexity. Gemma-2-2B-it is the primary Gemma baseline; Gemma-4-E2B reruns (3 epochs) show zero gain from gold $Z_t$ XML prefix conditioning (13.48 vs. 13.61).}
\label{fig:response_ppl}
\end{figure}

\begin{figure}[tp]
\centering
\includegraphics[width=0.88\linewidth]{figures/qwen_latent_training.pdf}
\caption{Qwen3-4B latent-predictor training dynamics. Mean accuracy keeps improving through epoch 4, while response-policy macro-F1 peaks earlier, suggesting a trade-off between broad head accuracy and decision-policy calibration.}
\label{fig:qwen_training}
\end{figure}

\begin{table}[tp]
\centering
\caption{Scenario-type distribution across train/val/test splits (episode-stratified 80/10/10).}
\label{tab:scenario_dist}
\footnotesize\setlength{\tabcolsep}{2.5pt}
\resizebox{\linewidth}{!}{%
\begin{tabular}{@{}lrrrrrr@{}}
\toprule
\textbf{Scenario} & \textbf{Train} & \textbf{\%} & \textbf{Val} & \textbf{\%} & \textbf{Test} & \textbf{\%} \\
\midrule
Secret extraction    & 1,395 & 22.6 & 113 & 16.5 & 158 & 17.9 \\
Rumor confrontation  & 1,091 & 17.7 & 128 & 18.8 & 173 & 19.6 \\
Apology repair       & 1,143 & 18.5 &  81 & 11.9 & 145 & 16.4 \\
Trust building       &   805 & 13.0 & 156 & 22.8 & 101 & 11.4 \\
Alliance negotiation &   677 & 11.0 &  81 & 11.9 &  82 &  9.3 \\
Threat escalation    &   569 &  9.2 &  49 &  7.2 &  83 &  9.4 \\
Deception detection  &   495 &  8.0 &  75 & 11.0 & 142 & 16.1 \\
\midrule
\textbf{Total}       & 6,175 & 100  & 683 & 100  & 884 & 100  \\
\bottomrule
\end{tabular}}
\end{table}

\begin{table}[tp]
\centering
\caption{Memory requirements per model. Training memory includes gradients and AdamW optimizer states; inference memory is weights only plus KV-cache at batch size 1. ``CPU'' means the model runs comfortably on a laptop with no GPU.}
\label{tab:hardware}
\footnotesize\setlength{\tabcolsep}{2pt}
\resizebox{\linewidth}{!}{%
\begin{tabular}{@{}lrrrrr@{}}
\toprule
\textbf{Model} & \textbf{Params} & \textbf{Precision} & \textbf{Disk} & \textbf{Train VRAM} & \textbf{Infer VRAM} \\
\midrule
GPT, PrefixGPT, MoE, Mamba (A) & 15--45M & fp32 & 0.06--0.18\,GB & 1--3\,GB & CPU \\
DistilBERT encoders (B)         & 66M      & fp32 & 0.26\,GB & 3--4\,GB & CPU \\
TinyLlama + LoRA (C)            & 1.1B     & 4-bit & 0.6\,GB  & 8--12\,GB & 4--6\,GB \\
Qwen3-4B + QLoRA (D)            & 4.0B     & 4-bit & 2.1\,GB  & 18--24\,GB& 8--12\,GB \\
Gemma\,4\,E2B + QLoRA (C)        & 5.1B     & 4-bit & 2.8\,GB  & 19--24\,GB& 12--16\,GB \\
\bottomrule
\end{tabular}}
\end{table}

\begin{table}[tp]
\centering
\caption{Multi-seed VAD placebo conditioning (TinyLlama-1.1B + LoRA, 3 epochs). All six runs converge to the same tight PPL range; real VAD provides no advantage over random VAD.}
\label{tab:placebo}
\footnotesize\setlength{\tabcolsep}{3pt}
\begin{tabular}{@{}lrrr@{}}
\toprule
\textbf{VAD Condition} & \textbf{Seed 42} & \textbf{Seed 43} & \textbf{Seed 44} \\
\midrule
Real OCEAN+VAD   & 2.90 & 2.86 & 2.88 \\
Random VAD       & 2.88 & 2.88 & 2.90 \\
\bottomrule
\end{tabular}
\end{table}

\section{Joint vs. Separate Training}
\label{app:joint}

\begin{table}[tp]
\centering
\caption{Joint vs. separate training per-head comparison (selected heads: those with largest $\Delta\kappa$). Full 28-head table in \texttt{eval\_results/ablation\_joint\_vs\_separate.json}.}
\label{tab:joint_vs_separate}
\scriptsize\setlength{\tabcolsep}{2pt}
\resizebox{\linewidth}{!}{%
\begin{tabular}{@{}lrrrrr@{}}
\toprule
\textbf{Head} & \textbf{Separate Acc} & \textbf{Joint Acc} & \textbf{Separate $\kappa$} & \textbf{Joint $\kappa$} & \textbf{$\Delta\kappa$} \\
\midrule
secrecy\_pressure  & 0.531 & 0.652 & 0.132 & 0.253 & \textbf{+0.121} \\
dominance\_delta   & 0.499 & 0.574 & 0.254 & 0.355 & \textbf{+0.101} \\
control            & 0.691 & 0.725 & 0.505 & 0.557 & +0.053 \\
threat             & 0.715 & 0.757 & 0.540 & 0.579 & +0.039 \\
risk\_type         & 0.861 & 0.870 & 0.358 & 0.403 & +0.045 \\
\midrule
response\_policy   & 0.698 & 0.684 & 0.607 & 0.581 & $-$0.026 \\
reveal\_decision   & 0.710 & 0.612 & 0.361 & 0.231 & $-$0.130 \\
familiarity\_level & 0.527 & 0.434 & 0.290 & 0.159 & $-$0.132 \\
respect\_level     & 0.600 & 0.508 & 0.462 & 0.341 & $-$0.122 \\
face\_pressure     & 0.879 & 0.865 & 0.216 & 0.110 & $-$0.106 \\
\midrule
\textbf{Mean (28 heads)} & 0.683 & 0.674 & 0.438 & 0.414 & $-$0.024 \\
\bottomrule
\end{tabular}}
\end{table}

Both models use the same Qwen3-4B backbone and were trained on identical data.
Joint training slightly degrades mean latent accuracy ($-$0.009) and mean $\kappa$ ($-$0.024) compared to the separate-stage pipeline, consistent with the well-known challenge of multi-task interference.
However, several heads improve under joint training: secrecy\_pressure ($\Delta\kappa{=}+0.121$) and dominance\_delta ($\Delta\kappa{=}+0.101$) benefit from the shared representation with the language modelling objective.
Consistency violations drop from 2 in the separate model to 0 for the critical \texttt{high\_secrecy\_full\_reveal} rule, suggesting the joint consistency loss ($\lambda{=}0.5$) provides a modest regularisation benefit.
The response generation perplexity is identical between the two settings (PPL~1.04 for both), indicating that joint training does not harm generation quality.
We treat this as a neutral-to-slightly-negative result: joint training is not harmful, but the consistency-loss benefits are too small to justify the additional architectural complexity at this data scale.

\section{Gold $Z_t$ Per-Field Response Difficulty}
\label{app:gold_zt_bridge}

To quantify how much each social-state field affects generation difficulty, we evaluate the Qwen3-4B response generator with gold $Z_t$ conditioning and measure per-field PPL variation across label values (Table~\ref{tab:gold_zt}).
If gold $Z_t$ values for a field produce systematically different PPLs, that field is response-relevant rather than merely predictable from context.

\begin{table}[tp]
\centering
\caption{Gold $Z_t$ per-field response PPL variation ($n=500$ examples). Effect size is Cohen's $d$ across label values.}
\label{tab:gold_zt}
\footnotesize
\resizebox{\linewidth}{!}{%
\begin{tabular}{@{}lrrr@{}}
\toprule
\textbf{Field} & \textbf{Min PPL} & \textbf{Max PPL} & \textbf{Effect $d$} \\
\midrule
reveal\_decision    & 7.23 (partial) & 11.73 (full) & 1.02 \\
response\_policy    & 7.34 (clarify) & 9.02 (partial) & 0.57 \\
repair\_strategy    & 8.08 (redirect) & 10.67 (apologize) & 0.59 \\
secrecy\_pressure   & 6.72 (low)    & 8.65 (medium) & 0.44 \\
player\_knowledge   & 7.84 (partial) & 9.56 (unaware) & 0.39 \\
trust\_delta        & 7.93 (neutral) & 10.07 (large+) & 0.48 \\
respect\_delta      & 8.07 (decrease)& 10.20 (large-) & 0.48 \\
dominance\_delta    & 7.25 (large+)  & 8.88 (large-) & 0.37 \\
\bottomrule
\end{tabular}}
\end{table}

Reveal decision has the largest causal effect on generation difficulty: generating a ``full reveal'' response is substantially harder than a ``hint'' or ``partial'' response.
Response policy shows moderate variation ($d{=}0.57$), with ``clarify'' being easiest and ``threaten'' hardest.
Relational deltas (trust, respect, dominance) show smaller effects, consistent with their poor latent predictability.

\section{Per-Head Evaluation Details}
\label{app:perhead}

\begin{figure}[tp]
\centering
\includegraphics[width=0.88\linewidth]{figures/latent_head_heatmap.pdf}
\caption{Per-head latent accuracy and true Cohen's $\kappa$ (Qwen3-4B). Several relational-delta heads are weak, but respect\_delta and trust\_delta retain measurable signal; affective arousal and valence are the easiest.}
\label{fig:latent_heatmap}
\end{figure}

\begin{table}[tp]
\centering
\caption{Per-head latent state prediction metrics (Qwen3-4B), sorted by true Cohen's $\kappa$ (descending). All $\kappa$ values are computed from confusion matrices.}
\label{tab:perhead}
\scriptsize\setlength{\tabcolsep}{1.5pt}
\begin{tabular}{@{}lrrrrr@{}}
\toprule
\textbf{Head} & \textbf{Group} & \textbf{\#Cls} & \textbf{Acc $\uparrow$} & \textbf{Macro-F1} & \textbf{Cohen's $\kappa$} \\
\midrule
arousal           & $A_t$ & 3  & 0.836 & 0.725 & 0.692 \\
valence           & $A_t$ & 3  & 0.818 & 0.752 & 0.688 \\
threat            & $A_t$ & 3  & 0.795 & 0.690 & 0.641 \\
response\_policy  & $D_t$ & 10 & 0.716 & 0.622 & 0.632 \\
duty\_pressure    & $N_t$ & 3  & 0.816 & 0.814 & 0.629 \\
trust\_level      & $R_t$ & 5  & 0.750 & 0.591 & 0.617 \\
player\_intent    & $M_t$ & 9  & 0.753 & 0.373 & 0.596 \\
respect\_level    & $R_t$ & 5  & 0.696 & 0.558 & 0.594 \\
respect\_delta    & $R_t$ & 5  & 0.714 & 0.514 & 0.584 \\
control           & $A_t$ & 3  & 0.725 & 0.689 & 0.561 \\
affection\_level  & $R_t$ & 5  & 0.701 & 0.645 & 0.552 \\
obligation\_level & $R_t$ & 5  & 0.666 & 0.449 & 0.508 \\
tone              & $C_t$ & 6  & 0.622 & 0.510 & 0.490 \\
player\_knowledge & $M_t$ & 4  & 0.638 & 0.611 & 0.487 \\
trust\_delta      & $R_t$ & 5  & 0.626 & 0.516 & 0.484 \\
repair\_strategy  & $D_t$ & 5  & 0.625 & 0.481 & 0.471 \\
risk\_type        & $C_t$ & 5  & 0.876 & 0.300 & 0.448 \\
obligation\_delta & $R_t$ & 5  & 0.679 & 0.349 & 0.443 \\
reveal\_decision  & $D_t$ & 4  & 0.687 & 0.656 & 0.420 \\
affection\_delta  & $R_t$ & 5  & 0.532 & 0.429 & 0.391 \\
familiarity\_level& $R_t$ & 5  & 0.578 & 0.522 & 0.389 \\
dominance\_level  & $R_t$ & 5  & 0.668 & 0.412 & 0.380 \\
player\_credibility & $M_t$ & 3 & 0.786 & 0.579 & 0.376 \\
dominance\_delta  & $R_t$ & 5  & 0.568 & 0.385 & 0.368 \\
value\_conflict   & $N_t$ & 3  & 0.729 & 0.557 & 0.344 \\
face\_pressure    & $N_t$ & 3  & 0.776 & 0.527 & 0.307 \\
familiarity\_delta& $R_t$ & 5  & 0.448 & 0.428 & 0.296 \\
secrecy\_pressure & $N_t$ & 3  & 0.587 & 0.491 & 0.176 \\
\bottomrule
\end{tabular}
\end{table}

\section{Sample Generations}
\label{app:samples}

\subsection{Data Examples}
\label{app:data_examples}

Table~\ref{tab:data_examples} shows two representative turns from the test split, with the full dialogue context, scene setup, and gold $Z_t$ labels.
These illustrate the input that the latent predictor receives (concatenated scene, stance history, and player utterance) and the 29-field label vector it must predict.

\begin{table*}[tp]
\centering
\caption{Representative test-set turns with full context and gold $Z_t$ labels.}
\label{tab:data_examples}
\footnotesize
\begin{tabular}{@{}p{3.8cm}p{11.5cm}@{}}
\toprule
\textbf{Field} & \textbf{Example 1: Apology Repair (turn 3)} \\
\midrule
Episode & 005d03da, turn 3, scenario=apology\_repair \\
Scene & Temple courtyard; NPC=priest; goals=protect\_flock, maintain\_authority; secrets=hidden\_artifact, church\_corruption \\
Player utterance & ``I fear my eyes see more than dust, Father. I've seen the enemy within---the whispers in the dark, the hands that reach for power.'' \\
\midrule
\textbf{Labels} & \\
dialogue\_act & [confess, probe] \\
tone & confrontational \\
risk\_type & secret-risk \\
valence / arousal / threat / control & negative / high / medium / medium \\
player\_intent / knowledge / credibility & probe / partial / medium \\
affection\_level / delta & VL / 0 \\
respect\_level / delta & VL / 0 \\
dominance\_level / delta & H / ++ \\
familiarity\_level / delta & N / 0 \\
trust\_level / delta & VL / 0 \\
obligation\_level / delta & H / ++ \\
duty / secrecy / face pressure & high / high / medium \\
value\_conflict & strong \\
response\_policy / reveal\_decision / repair\_strategy & challenge / hint / redirect \\
\bottomrule
\end{tabular}
\end{table*}

\subsection{Prediction Examples}
\label{app:prediction_examples}

Table~\ref{tab:prediction_examples} shows the Qwen3-4B latent predictor's predictions on four held-out test turns, with gold labels and correctness indicated.
These examples illustrate both the model's strengths (\texttt{value\_conflict} and \texttt{secrecy\_pressure} are often correct) and weaknesses (\texttt{response\_policy} and \texttt{reveal\_decision} are frequently confused).

\begin{table*}[tp]
\centering
\caption{Latent predictor predictions on held-out test turns. \cmark=correct, \xmark=incorrect.}
\label{tab:prediction_examples}
\footnotesize\setlength{\tabcolsep}{3pt}
\begin{tabular}{@{}lllll@{}}
\toprule
\textbf{Scenario} & \textbf{Field} & \textbf{Gold} & \textbf{Predicted} & \\
\midrule
Alliance negotiation & response\_policy  & threaten & deflect  & \xmark \\
                   & reveal\_decision   & hint     & none     & \xmark \\
                   & secrecy\_pressure  & high     & high     & \cmark \\
                   & value\_conflict    & strong   & strong   & \cmark \\
\midrule
Apology repair     & response\_policy  & clarify  & deflect  & \xmark \\
                   & reveal\_decision   & hint     & hint     & \cmark \\
                   & secrecy\_pressure  & medium   & high     & \xmark \\
                   & value\_conflict    & mild     & mild     & \cmark \\
\midrule
Deception detection& response\_policy  & deflect  & challenge& \xmark \\
                   & reveal\_decision   & hint     & hint     & \cmark \\
                   & secrecy\_pressure  & medium   & high     & \xmark \\
                   & value\_conflict    & mild     & mild     & \cmark \\
\midrule
Rumor confrontation& response\_policy  & deflect  & partial  & \xmark \\
                   & reveal\_decision   & none     & hint     & \xmark \\
                   & secrecy\_pressure  & high     & high     & \cmark \\
                   & value\_conflict    & strong   & strong   & \cmark \\
\bottomrule
\end{tabular}
\end{table*}

\subsection{Generated Response Examples}

Table~\ref{tab:samples} shows representative generated responses from the Qwen3-4B response model.

\begin{table*}[tp]
\centering
\caption{Sample NPC responses. Latent state context is abbreviated.}
\label{tab:samples}
\footnotesize\setlength{\tabcolsep}{2pt}
\begin{tabular}{@{}p{3.8cm}p{5.2cm}p{5.2cm}r@{}}
\toprule
\textbf{Latent Context} & \textbf{Gold Response} & \textbf{Generated Response} & \textbf{ROUGE-L} \\
\midrule
policy=negotiate, reveal=none, secrecy=medium &
``I can mend your blade, but I won't do it for nothing. You'll have to trade something of value...'' &
``Okay, I can fix your sword, but I need a fair price for my time and materials. How much do you think you can offer?...'' &
0.103 \\
\midrule
policy=clarify, reveal=hint, secrecy=high &
``I've heard whispers, not from the council, but from those who know the iron in their bones...'' &
``The Iron Legion is a band of warriors from the east... I've heard whispers that they're looking for a blade of ancient power...'' &
0.104 \\
\bottomrule
\end{tabular}
\end{table*}

\section{Data Generation Pipeline and Provenance}
\label{app:datagen}

\subsection{Generation Pipeline Architecture}

Figure~\ref{fig:data_flow} illustrates the full pipeline.
The system comprises nine stages executed in sequence:

\begin{enumerate}
    \item \textbf{ScenarioBank:} Loads 7 YAML scenario definitions (secret extraction, trust building, alliance negotiation, threat escalation, deception detection, rumor confrontation, apology repair) plus a persistent world context (``Oakhaven Siege'').
    Each scenario specifies NPC archetypes, initial relational stances, conflict drivers, and domain-specific constraints on legal $Z_t$ transitions.
    \item \textbf{StateInit:} Samples an initial social state $Z_0$ for each NPC in the episode, respecting the scenario's constraints.
    \item \textbf{EpisodePlanner:} Generates a high-level plan for the conversation: player goal, NPC goal, expected number of turns (5--10), and which social-state dimensions should evolve.
    \item \textbf{TurnGenerator:} For each turn, constructs a prompt containing the scene description, conversation history, and current $Z_t$, then queries the teacher LLM to produce the next utterance.
    Supports multiple teacher backends (local HuggingFace models, Azure OpenAI, OpenRouter) with parallel worker threads for throughput.
    \item \textbf{Labeler:} A specialised 10-step teacher LLM pipeline that annotates each generated turn with the full 29-dimension $Z_t$ schema.
    Each labelling step focuses on a subset of dimensions (e.g., first label dialogue act and tone, then affect, then relational stance, etc.) to reduce annotation errors.
    \item \textbf{CounterfactualAugmenter:} Produces 3 counterfactual variants per episode by flipping selected $Z_t$ dimensions (e.g., reversing secrecy pressure, swapping response policy) and regenerating the affected turns.
    \item \textbf{Validator:} Checks each turn against the schema's consistency rules (e.g., full reveal requires low secrecy pressure; answer policy requires a reveal decision).
    Invalid turns are discarded or re-queued for regeneration.
    \item \textbf{Packager:} Converts validated turns into the SFT format (\texttt{input} $\rightarrow$ \texttt{target}) and the heads-supervision format (\texttt{context} $\rightarrow$ \texttt{labels}).
    \item \textbf{Splitter:} Performs an 80/10/10 episode-stratified split into train, validation, and test sets, ensuring no episode appears in more than one split.
\end{enumerate}

The pipeline is configurable via \texttt{configs/data\_gen.yaml}: target 5,000 episodes, 5--10 turns per episode, 3 counterfactuals each, with multi-provider teacher pooling for parallel generation.
A dry-run mode (\texttt{--dry-run}) uses mock LLM calls for rapid testing without API costs.

\subsection{Data Sources and Sizes}

Table~\ref{tab:datasources} documents the provenance of every dataset used in training and evaluation.

\begin{table*}[tp]
\centering
\caption{Data sources, sizes, and formats. ``Synthetic'' means generated by a teacher LLM; ``Real'' means from published academic corpora.}
\label{tab:datasources}
\footnotesize\setlength{\tabcolsep}{3pt}
\begin{tabular}{@{}llrlp{0.42\linewidth}@{}}
\toprule
\textbf{Track} & \textbf{Dataset} & \textbf{Size} & \textbf{Type} & \textbf{Origin} \\
\midrule
A & SLM dialogue text & 16,905 lines & Synthetic & Teacher LLM $\times$ 7 scenario templates $\times$ 414 NPC profiles \\
A & Validation text & 891 lines & Synthetic & Held-out 5\% split \\
\midrule
B & Personality (OCEAN) & 2,221 essays & \textbf{Real} & \texttt{jingjietan/essays-big5} (HuggingFace); student essays with Big Five scores \\
B & Affect (VAD) & 9,056 samples & \textbf{Real} & EmoBank (GitHub/JULIELab); human-rated valence-arousal-dominance annotations \\
\midrule
C+D & SFT dialogue & 7,742 turns & Synthetic & Teacher LLM; 6,175 train / 683 val / 884 test; 736 episodes, 7 scenarios \\
C+D & Heads labels & 7,742 turns & Synthetic & Teacher LLM; same turns as SFT, annotated with full 29-dim $Z_t$ schema \\
C & Gemma dialogue & 2,298 turns & Synthetic & Structured JSONL with NPC profiles; 2,183 train / 115 val \\
\midrule
--- & NPC profiles & 414 entries & Synthetic & Generated from scenario bank (apothecary, guard, merchant, spy, scholar, etc.) \\
\bottomrule
\end{tabular}
\end{table*}

\subsection{External Corpora}

Four published dialogue datasets and two emotion/personality corpora were downloaded and merged into a unified format:

\begin{itemize}
    \item \textbf{PersonaChat}~\cite{zhang2018personalizing}: persona-conditioned chit-chat (ACL 2018); extracted dialogue and personality profiles.
    \item \textbf{EmpatheticDialogues}~\cite{rashkin2019empathetic}: emotion-grounded conversations (ACL 2019); extracted dialogue and per-utterance emotion labels.
    \item \textbf{LIGHT}~\cite{urbanek2019learning}: character-and-setting-conditioned fantasy dialogue (EMNLP 2019); extracted dialogue and personality profiles.
    \item \textbf{PIPPA}~\cite{gosmar2022pippa}: open-domain personalized conversational dataset (arXiv 2022); extracted dialogue and personality profiles.
    \item \textbf{Essays-Big5}\footnote{\url{https://huggingface.co/datasets/jingjietan/essays-big5}}: 1,578 student essays with self-reported Big Five OCEAN scores; hosted on HuggingFace.
    \item \textbf{EmoBank}~\cite{buechel2017emobank}: 10,548 sentences with human-rated valence, arousal, and dominance annotations (EACL 2017); sourced from GitHub (JULIELab).
\end{itemize}

These were merged into four unified files:
\begin{itemize}
    \item \texttt{merged\_dialogue.jsonl}: 70,764 dialogue turns combining PersonaChat, EmpatheticDialogues, LIGHT, and PIPPA.
    \item \texttt{merged\_dialogue.txt}: 2.7M characters of plain-text dialogue for SLM pretraining.
    \item \texttt{merged\_personality.jsonl}: 51,231 personality profiles from PersonaChat, LIGHT, and PIPPA.
    \item \texttt{merged\_affect.csv}: 19,534 affect-labelled utterances from EmpatheticDialogues and EmoBank.
\end{itemize}

\textbf{Training subset:} From these merged corpora, we sampled 2,221 personality-labelled examples and 9,056 affect-labelled examples for encoder training (Track~B).
The SLM dialogue corpus (16,905 lines; Track~A) was drawn from \texttt{merged\_dialogue.txt}.
The full merged corpora (70K+ turns) will be released upon acceptance for full reproducibility.

\textbf{Distinction from synthetic data:} The personality (essays-big5) and affect (EmoBank, EmpatheticDialogues) labels are \textit{real human annotations} from published academic corpora.
The dialogue text in the external corpora is also real (crowdworker-generated), in contrast to the synthetic teacher-LLM-generated SFT data used for Tracks~C and~D.

\subsection{Label Class Distributions}

Table~\ref{tab:label_dist} reports the class distributions for the six most actionable heads (all single-label).
The data exhibits moderate class imbalance: \texttt{reveal\_decision} is dominated by \textit{hint} and \textit{none} (96.6\% combined), \texttt{response\_policy} is dominated by \textit{deflect} and \textit{soothe} (59.7\% combined), and \texttt{secrecy\_pressure} is heavily skewed toward \textit{high} and \textit{medium} (97.2\%).
This imbalance motivates the focal-loss and class-weighting strategies used in Track~D training.
Stance-level fields contain label-noise artifacts (e.g., ``H(+)'', ``VH(--)'') from early-generation labelling errors; these were cleaned in the final training split but remain in the raw distribution counts.

\begin{table*}[tp]
\centering
\caption{Class distributions for key single-label heads (all 7,742 turns).}
\label{tab:label_dist}
\footnotesize\setlength{\tabcolsep}{2.5pt}
\begin{tabular}{@{}llrrrrrr@{}}
\toprule
\textbf{Head} & \textbf{Most freq.} & \textbf{\%} & \textbf{2nd} & \textbf{\%} & \textbf{3rd} & \textbf{\%} & \textbf{\#Classes} \\
\midrule
response\_policy  & deflect   & 33.6 & soothe    & 26.1 & threaten   & 11.5 & 13 \\
reveal\_decision   & hint      & 54.9 & none      & 40.6 & partial    &  3.6 & 5  \\
secrecy\_pressure  & high      & 54.1 & medium    & 43.0 & (low/empty)&  2.9 & 3  \\
duty\_pressure     & high      & 56.6 & medium    & 43.1 & low        &  0.2 & 3  \\
player\_intent     & probe     & 58.1 & intimidate& 12.0 & seek-info  &  8.5 & 9  \\
valence           & negative  & 58.9 & positive  & 22.2 & neutral    & 18.9 & 3  \\
\bottomrule
\end{tabular}
\end{table*}

\subsection{SLM Training Corpus Distribution}

Table~\ref{tab:slm_dist} summarises the Track~A SLM training corpus (2,183 structured records from \texttt{from\_gen\_train.jsonl}).
This corpus is drawn from the same generated episodes as Tracks~C and~D, but packaged as plain-text dialogue lines for language-model pretraining.

\begin{table*}[tp]
\centering
\caption{SLM training corpus distribution (2,183 records, 396 unique NPCs).}
\label{tab:slm_dist}
\footnotesize\setlength{\tabcolsep}{2.5pt}
\begin{tabular}{@{}lrlrlr@{}}
\toprule
\textbf{Value} & \textbf{Count} & \textbf{\%} & \textbf{Value} & \textbf{Count} & \textbf{\%} \\
\midrule
\multicolumn{6}{l}{\textit{Scenario type}} \\
Secret extraction    & 439 & 20.1 & Rumor confrontation & 401 & 18.4 \\
Apology repair       & 376 & 17.2 & Trust building      & 317 & 14.5 \\
Alliance negotiation & 253 & 11.6 & Deception detection & 227 & 10.4 \\
Threat escalation    & 170 &  7.8 & & & \\
\midrule
\multicolumn{6}{l}{\textit{Response policy}} \\
Deflect   & 811 & 37.1 & Soothe     & 441 & 20.2 \\
Challenge & 273 & 12.5 & Threaten & 197 &  9.0 \\
\midrule
\multicolumn{6}{l}{\textit{Valence / Arousal}} \\
Negative  & 1,359 & 62.3 & High     & 1,085 & 49.7 \\
Neutral   &   419 & 19.2 & Medium   & 1,062 & 48.6 \\
Positive  &   405 & 18.6 & Low      &    36 &  1.6 \\
\midrule
\multicolumn{6}{l}{\textit{Lengths (words)}} \\
Mean response length & \multicolumn{2}{r}{29.8} & Mean context length & \multicolumn{2}{r}{141.1} \\
\bottomrule
\end{tabular}
\end{table*}

\subsection{Limitation: Synthetic Labels}
We note that all dialogue data and social-state annotations are LLM-generated, not human-annotated.
This introduces label noise, particularly for subjective dimensions like relational stance and tone.
Results should be interpreted with this in mind; human evaluation of label quality is planned for future work.

\section{Real-World Use Cases and Practical Accomplishments}
\label{app:usecases}

See Supplementary Material, Appendix~F for deployment scenarios, inference costs, and a capabilities summary table.

\section{System Data Flow}
\label{app:dataflow}

See Supplementary Material, Appendix~G for full input$\rightarrow$model$\rightarrow$output specifications for each track, including tokenisation details, loss functions, and checkpoint paths.

\section{Training Pipeline Infrastructure \& Code Modifications}
\label{app:infrastructure}

See Supplementary Material, Appendix~H for Slurm job management scripts, Qwen3 training configurations, the evaluation framework, and a table documenting the 16 code modifications required for reproducibility.


\section{Social-State JEPA Details}
\label{app:jepa}

We evaluated a Social-State JEPA auxiliary objective~\cite{lecun2022path,assran2023self} in which a predictor network is trained to predict future $Z_{t+1}$ representations in embedding space, with a stop-gradient on the target encoder.
The JEPA objective was tested on four backbones (Qwen3-4B, Gemma-4-E2B, from-scratch GPT-16M, from-scratch Mamba-15M) and compared against a shuffled-future placebo on Qwen3-4B and Mamba-15M.

\begin{table}[tp]
\centering
\caption{JEPA per-backbone results. Positive $\Delta$ means JEPA improves over base. Shuffled-future placebo partially replicates gains, indicating temporal structure is not the sole driver.}
\label{tab:jepa_appendix}
\footnotesize\setlength{\tabcolsep}{2pt}
\resizebox{\linewidth}{!}{%
\begin{tabular}{@{}lrrrrrrr@{}}
\toprule
\textbf{Backbone} & \textbf{Base $\kappa$} & \textbf{JEPA $\kappa$} & \textbf{$\Delta\kappa$} & \textbf{Shuf $\kappa$} & \textbf{Heads $\uparrow$} & \textbf{Heads $\downarrow$} & \textbf{Sig $\downarrow$} \\
\midrule
Qwen3-4B     & 0.441 & 0.468 & +0.027 & 0.441 & 12 & 16 & 10 \\
GPT-16M      & 0.370 & 0.387 & +0.017 & ---   & 14 & 14 & 9 \\
Mamba-15M    & 0.148 & 0.159 & +0.012 & 0.127 & 12 & 16 & 10 \\
Gemma-4-E2B  & ---   & 0.546 & ---    & ---   & --- & --- & --- \\
\bottomrule
\end{tabular}}
\end{table}

Across all backbones with base comparisons, the JEPA objective produces small mean $\kappa$ improvements ($+0.012$ to $+0.027$) but more heads degrade than improve (12--16 heads degrade vs.\ 12--14 improve), and 9--10 heads are significantly worse under JEPA.
On Qwen3-4B, the shuffled-future placebo ($\kappa{=}0.441$) nearly matches the base model ($\kappa{=}0.441$) but does not replicate the full JEPA gain ($\kappa{=}0.468$), leaving open the possibility of a small temporal benefit that is not statistically robust.
On Mamba-15M, the shuffled placebo ($\kappa{=}0.127$) is actually \textit{worse} than both base ($\kappa{=}0.148$) and JEPA ($\kappa{=}0.159$).
Gemma-4-E2B JEPA evaluation (no base comparison) achieves mean accuracy 0.546.
We report the JEPA experiment as a mixed-to-null result: small aggregate gains are offset by significant per-head degradation, the shuffled placebo partially replicates the effect, and no single actionable head (response\_policy, reveal\_decision) shows a reliable improvement.
The negative finding is consistent with the broader observation that some per-turn deltas are weakly recoverable for supervised models as well, suggesting that single-turn social-state prediction has fundamental limits that architectural innovations like JEPA do not overcome without richer temporal context.

\section{Human Label Audit}
\label{app:human_audit}

\subsection{Audit Protocol}

To assess the audibility of the synthetic social-state labels, we conducted a stratified audit of 150 natural, non-counterfactual test-set turns (approximately 21 per scenario type) drawn from a cleaned 356-turn audit packet.
Five human annotators (recruited via Prolific) independently labelled each turn on a subset of eight heads selected for their downstream importance:
\texttt{valence}, \texttt{arousal}, \texttt{secrecy\_pressure}, \texttt{reveal\_decision}, \texttt{response\_policy}, \texttt{repair\_strategy}, \texttt{trust\_level}, and \texttt{familiarity\_level}.
Annotators were shown the scene description, conversation history, player utterance, and NPC response, and selected labels from dropdown menus matching the schema vocabulary.
No annotator saw model predictions during judgement.
We also ran a model-based validator (Gemma-4-31B-IT, zero-shot) on the same 150 turns as a third judge; this is reported separately from human--human agreement and is not treated as a replacement for human audit.

\subsection{Quality-Control Screening}

All five annotators were flagged by post-hoc quality-control criteria (Table~\ref{tab:audit_qc}).
Two annotators (H1, H2) completed the task at extreme speed (median 41--44\,s per turn, with 62--80\% of turns at the 50\,s minimum timer), suggesting speed-running rather than genuine engagement.
A third (H4) selected the same label for \texttt{familiarity\_level} on all 150 turns and the same label for \texttt{repair\_strategy} on 147/150 turns, indicating default-option bias.
All annotators left notes empty on $>$91\% of turns.
We therefore report agreement statistics for all annotators but flag that the human numbers should be interpreted as a \textbf{lower bound} on what trained, calibrated annotators could achieve.

\begin{table}[tp]
\centering
\caption{Annotator quality-control summary. Min-time ratio = fraction of turns at the 50\,s floor timer. Empty notes = fraction of turns with no free-text note.}
\label{tab:audit_qc}
\scriptsize
\resizebox{\linewidth}{!}{%
\begin{tabular}{@{}lrrrrl@{}}
\toprule
\textbf{ID} & \textbf{Turns} & \textbf{Med.\ time (s)} & \textbf{Min-time \%} & \textbf{Empty notes \%} & \textbf{Primary QC flag} \\
\midrule
H1 & 150 & 41 & 80\% & 100\% & Speed-runner \\
H2 & 160 & 44 & 62\% & 99\% & Speed-runner \\
H3 & 150 & 60 & 41\% & 91\% & Rushed \\
H4 & 150 & 55 & 41\% & 100\% & Default-option bias \\
H5 & 82 & 113 & 4\% & 100\% & Timer abuse \\
\bottomrule
\end{tabular}}
\end{table}

\subsection{Model-Based Validator vs.\ Teacher (Internal Consistency)}

Table~\ref{tab:ai_teacher} reports agreement between the model-based validator and the teacher labels.
Because both are LLM-generated, this comparison measures \textbf{internal consistency} of the schema in model space rather than human-grounded validity.
The model-based validator achieves $\kappa{=}0.40$ on \texttt{response\_policy} and $\kappa{=}0.32$ on \texttt{valence} with the teacher, which is consistent with the interpretation that these heads carry recoverable signal across independently generated model outputs.
Heads with low model--teacher agreement---\texttt{trust\_level} ($\kappa{=}0.02$), \texttt{familiarity\_level} ($\kappa{=}0.02$), and \texttt{arousal} ($\kappa{=}0.05$)---indicate dimensions where even models disagree, suggesting inherent ambiguity or underspecification in the labelling guidelines.

\begin{table}[tp]
\centering
\caption{Model-based validator vs.\ teacher labels ($n{=}150$ turns). Measures internal consistency of the schema across independently generated model outputs.}
\label{tab:ai_teacher}
\footnotesize\setlength{\tabcolsep}{3pt}
\begin{tabular}{@{}lrr@{}}
\toprule
\textbf{Head} & \textbf{M--Teacher Acc} & \textbf{M--Teacher $\kappa$} \\
\midrule
valence              & 0.527 & 0.316 \\
arousal              & 0.367 & 0.045 \\
secrecy\_pressure    & 0.540 & 0.147 \\
reveal\_decision     & 0.520 & 0.148 \\
response\_policy     & 0.507 & 0.397 \\
repair\_strategy     & 0.413 & 0.136 \\
trust\_level         & 0.207 & 0.022 \\
familiarity\_level   & 0.367 & 0.019 \\
\midrule
\textbf{Mean}        & \textbf{0.431} & \textbf{0.154} \\
\bottomrule
\end{tabular}
\end{table}

\subsection{Human--Teacher Agreement}

Table~\ref{tab:human_teacher} reports per-annotator agreement with the teacher labels.
Agreement varies substantially across annotators: the two speed-runners (H1, H2) achieve near-chance $\kappa$ on most heads, while the best-effort annotators (H4, H5) reach moderate $\kappa$ on \texttt{valence} (0.27--0.28), \texttt{response\_policy} (0.31--0.37), and \texttt{reveal\_decision} (0.18--0.19).
Even the best human annotator, however, achieves $\kappa{<}0.10$ on \texttt{familiarity\_level}, \texttt{repair\_strategy}, and \texttt{arousal}, which is consistent with the interpretation that these heads resist reliable human annotation under the current protocol.

\begin{table}[tp]
\centering
\caption{Per-annotator human--teacher agreement ($\kappa$). H1--H4 annotated 150 turns; H5 annotated 82. Best $\kappa$ per head is \textbf{bolded}.}
\label{tab:human_teacher}
\footnotesize\setlength{\tabcolsep}{2pt}
\resizebox{\linewidth}{!}{%
\begin{tabular}{@{}lrrrrr@{}}
\toprule
\textbf{Head} & \textbf{H1} & \textbf{H2} & \textbf{H3} & \textbf{H4} & \textbf{H5} \\
\midrule
valence              & 0.276 & $-$0.024 & 0.134 & \textbf{0.271} & 0.280 \\
arousal              & 0.021 & $-$0.043 & 0.049 & 0.027 & \textbf{0.107} \\
secrecy\_pressure    & $-$0.013 & $-$0.003 & 0.165 & \textbf{0.197} & 0.158 \\
reveal\_decision     & $-$0.001 & $-$0.007 & 0.070 & \textbf{0.192} & 0.183 \\
response\_policy     & 0.322 & $-$0.043 & 0.204 & \textbf{0.374} & 0.309 \\
repair\_strategy     & \textbf{0.085} & 0.004 & $-$0.015 & 0.015 & 0.012 \\
trust\_level         & 0.076 & 0.000 & 0.070 & 0.094 & \textbf{0.148} \\
familiarity\_level   & 0.003 & $-$0.031 & $-$0.044 & 0.000 & \textbf{0.021} \\
\midrule
\textbf{Mean $\kappa$} & 0.096 & $-$0.018 & 0.079 & \textbf{0.146} & 0.152 \\
\bottomrule
\end{tabular}}
\end{table}

\subsection{Human--Human Inter-Annotator Agreement}

Table~\ref{tab:human_human} reports pairwise Cohen's $\kappa$ for all annotator pairs with $\geq$50 common turns.
Human--human agreement is near zero or negative for most heads and most pairs, consistent with the interpretation that the eight-head annotation task is difficult for uncalibrated annotators.
The single exception is the H4--H5 pair, which achieves substantial agreement on \texttt{valence} ($\kappa{=}0.79$), \texttt{reveal\_decision} ($\kappa{=}0.56$), and \texttt{trust\_level} ($\kappa{=}0.44$).
However, this pair also shows $\kappa{=}0.00$ on \texttt{familiarity\_level} (both default to ``N'') and $\kappa{=}-0.01$ on \texttt{repair\_strategy} (both default to ``none''), indicating that their agreement on other heads may partly reflect shared default-option bias rather than genuine convergence on the labelling criteria.
The two speed-runners (H1--H2) achieve $\kappa{=}0.009$ on average, essentially at chance.

\begin{table}[tp]
\centering
\caption{Pairwise human--human Cohen's $\kappa$ for annotator pairs with $\geq$50 common turns. H1--H4 share 150 turns; pairs involving H5 share 82. Pairs are ordered by mean $\kappa$ (descending).}
\label{tab:human_human}
\footnotesize\setlength{\tabcolsep}{1.5pt}
\resizebox{\linewidth}{!}{%
\begin{tabular}{@{}lrrrrrrrrrr@{}}
\toprule
 & \rotatebox{70}{H4--H5} & \rotatebox{70}{H1--H4} & \rotatebox{70}{H1--H5} & \rotatebox{70}{H1--H3} & \rotatebox{70}{H3--H5} & \rotatebox{70}{H3--H4} & \rotatebox{70}{H2--H4} & \rotatebox{70}{H2--H5} & \rotatebox{70}{H2--H3} & \rotatebox{70}{H1--H2} \\
\midrule
valence       & 0.79 & 0.44 & 0.40 & 0.37 & 0.30 & 0.42 & 0.02 & $-$0.09 & 0.03 & $-$0.06 \\
arousal       & 0.32 & 0.13 & 0.14 & 0.06 & 0.12 & 0.04 & 0.04 & 0.06 & $-$0.17 & 0.07 \\
secr.\ pres.  & 0.25 & $-$0.01 & $-$0.02 & 0.00 & 0.18 & 0.08 & $-$0.08 & $-$0.03 & $-$0.02 & $-$0.01 \\
reveal dec.   & 0.56 & $-$0.01 & 0.08 & 0.13 & 0.02 & 0.01 & $-$0.01 & $-$0.06 & $-$0.05 & 0.03 \\
resp.\ policy & 0.42 & 0.26 & 0.27 & 0.16 & 0.24 & 0.21 & $-$0.05 & $-$0.04 & 0.01 & 0.00 \\
repair strat. & $-$0.01 & $-$0.03 & $-$0.02 & 0.01 & 0.00 & $-$0.02 & 0.00 & 0.02 & $-$0.01 & $-$0.01 \\
trust level   & 0.44 & 0.31 & 0.17 & 0.12 & 0.13 & 0.10 & $-$0.01 & 0.02 & $-$0.02 & 0.05 \\
fam.\ level   & 0.00 & 0.00 & 0.02 & 0.06 & 0.01 & 0.00 & 0.00 & 0.01 & $-$0.01 & 0.01 \\
\midrule
\textbf{Mean} & \textbf{0.35} & \textbf{0.14} & \textbf{0.13} & \textbf{0.11} & \textbf{0.13} & \textbf{0.10} & \textbf{$-$0.01} & \textbf{$-$0.01} & \textbf{$-$0.03} & \textbf{0.01} \\
\bottomrule
\end{tabular}}
\end{table}

\subsection{Human--Model-Based Validator Agreement}

Table~\ref{tab:human_ai} reports agreement between each human annotator and the model-based validator.
The best-effort annotators (H4, H5) achieve moderate-to-substantial agreement with the model-based validator on \texttt{valence} ($\kappa{=}0.69$ and $0.76$), \texttt{response\_policy} ($\kappa{=}0.49$ and $0.53$), and \texttt{reveal\_decision} ($\kappa{=}0.44$ and $0.46$), suggesting that these heads capture signal that is consistent across human and model judgement.
In contrast, \texttt{familiarity\_level} and \texttt{repair\_strategy} show near-zero or negative human--model $\kappa$ for all annotators, which is consistent with the interpretation that these are the most ambiguous heads.

\begin{table}[tp]
\centering
\caption{Per-annotator human--model-based validator agreement ($\kappa$). H1--H4 share 150 turns with the model-based validator; H5 shares 82.}
\label{tab:human_ai}
\footnotesize\setlength{\tabcolsep}{2pt}
\resizebox{\linewidth}{!}{%
\begin{tabular}{@{}lrrrrr@{}}
\toprule
\textbf{Head} & \textbf{H1} & \textbf{H2} & \textbf{H3} & \textbf{H4} & \textbf{H5} \\
\midrule
valence              & 0.517 & $-$0.035 & 0.380 & \textbf{0.685} & 0.759 \\
arousal              & 0.171 & 0.023 & 0.028 & \textbf{0.411} & 0.218 \\
secrecy\_pressure    & 0.010 & 0.005 & 0.207 & 0.289 & \textbf{0.532} \\
reveal\_decision     & 0.069 & $-$0.019 & 0.043 & \textbf{0.440} & 0.455 \\
response\_policy     & 0.362 & $-$0.036 & 0.275 & \textbf{0.489} & 0.527 \\
repair\_strategy     & \textbf{0.198} & 0.002 & 0.049 & $-$0.021 & $-$0.093 \\
trust\_level         & 0.217 & $-$0.004 & 0.180 & \textbf{0.380} & 0.381 \\
familiarity\_level   & $-$0.013 & $-$0.039 & 0.111 & 0.000 & $-$0.030 \\
\midrule
\textbf{Mean $\kappa$} & 0.191 & $-$0.013 & 0.159 & \textbf{0.334} & 0.344 \\
\bottomrule
\end{tabular}}
\end{table}

\subsection{Consolidated Audit Summary}

Table~\ref{tab:audit_summary} consolidates the three agreement axes---human--teacher (best annotator), human--human (range across pairs), and model--teacher---into a single per-head summary.
Heads fall into three tiers:

\begin{enumerate}
    \item \textbf{Partially auditable:} \texttt{valence} and \texttt{response\_policy} achieve the highest agreement across all axes, with best-human--teacher $\kappa{\geq}0.27$, maximum HH $\kappa{\geq}0.40$, and model--teacher $\kappa{\geq}0.32$.
    \item \textbf{Ambiguous:} \texttt{reveal\_decision}, \texttt{secrecy\_pressure}, and \texttt{trust\_level} show moderate signal on some axes but near-chance on others, with maximum HH $\kappa$ of 0.56, 0.25, and 0.44 respectively but mean HH $\kappa$ near zero.
    \item \textbf{Non-auditable:} \texttt{familiarity\_level}, \texttt{repair\_strategy}, and \texttt{arousal} show near-zero or negative $\kappa$ on all axes, suggesting that these heads are either inherently ambiguous given the dialogue context or require substantially different annotation protocols.
\end{enumerate}

\begin{table}[tp]
\centering
\caption{Consolidated audit summary per head. Best HT $\kappa$ = highest human--teacher $\kappa$ across five annotators. HH $\kappa$ range = min--max across all pairs with $\geq$50 common turns. M-T $\kappa$ = model-based validator vs.\ teacher.}
\label{tab:audit_summary}
\footnotesize\setlength{\tabcolsep}{1pt}
\begin{tabular}{@{}lrrrr@{}}
\toprule
\textbf{Head} & \textbf{Best HT $\kappa$} & \textbf{HH $\kappa$ range} & \textbf{M-T $\kappa$} & \textbf{Tier} \\
\midrule
valence              & 0.280 & $-$0.09--0.79 & 0.316 & Partial \\
response\_policy     & 0.374 & $-$0.05--0.42 & 0.397 & Partial \\
reveal\_decision     & 0.192 & $-$0.06--0.56 & 0.148 & Ambiguous \\
secrecy\_pressure    & 0.197 & $-$0.08--0.25 & 0.147 & Ambiguous \\
trust\_level         & 0.148 & $-$0.02--0.44 & 0.022 & Ambiguous \\
arousal              & 0.107 & $-$0.17--0.32 & 0.045 & Non-audit. \\
repair\_strategy     & 0.085 & $-$0.03--0.02 & 0.136 & Non-audit. \\
familiarity\_level   & 0.021 & $-$0.01--0.06 & 0.019 & Non-audit. \\
\bottomrule
\end{tabular}
\end{table}

\subsection{Interpretation}

The human audit does not support treating the synthetic labels as human-validated gold annotations.
Human--human $\kappa$ is near zero for most heads and most annotator pairs, and all five annotators were flagged by post-hoc quality-control criteria for speed, default-option bias, or timer abuse.
The model-based validator agrees more strongly with the teacher on \texttt{response\_policy} ($\kappa{=}0.40$) and \texttt{valence} ($\kappa{=}0.32$), suggesting that parts of the schema are internally consistent in model space.
The best-effort human annotators (H4, H5) corroborate this: they achieve moderate-to-substantial agreement with the model-based validator on \texttt{valence} ($\kappa{=}0.69$--$0.76$) and \texttt{response\_policy} ($\kappa{=}0.49$--$0.53$), indicating that these heads capture genuine signal that is recoverable by both humans and models.
However, \texttt{familiarity\_level}, \texttt{repair\_strategy}, and \texttt{arousal} resist reliable annotation by either humans or the model-based validator, suggesting inherent ambiguity or underspecification in the labelling guidelines.

We therefore interpret $Z_t$ as useful teacher-generated supervision for controlled experiments, not as a validated model of human social judgement.
Future work should incorporate annotator training, calibration examples with gold labels, and adjudication rounds before using these labels as human-validated gold annotations.

\end{document}
